% ============================================================================
% CAPÍTULO 1 - INTRODUÇÃO E GUIA DE REPRODUÇÃO
% ============================================================================

\chapter{Introdução}

\section{Contexto e Motivação}

A \acrfull{ia} generativa tem transformado diversas áreas do conhecimento, e o setor de saúde desponta como um dos campos com maior potencial de impacto. A capacidade dos \acrfull{llm} de processar, interpretar e gerar texto em linguagem natural abre oportunidades significativas para auxiliar profissionais de saúde em tarefas como consulta a protocolos clínicos, suporte à tomada de decisão e acesso rápido a informações de pacientes.

No entanto, modelos generalistas de linguagem --- apesar de seu amplo conhecimento --- apresentam limitações importantes no domínio médico: podem gerar informações imprecisas, desatualizadas ou descontextualizadas. Essas limitações tornam necessário o uso de técnicas especializadas como \textbf{fine-tuning} (ajuste fino) para aprimorar o conhecimento médico do modelo, \textbf{RAG} (\textit{Retrieval-Augmented Generation}) para fundamentar respostas em protocolos clínicos reais, e \textbf{guardrails} de segurança para garantir que o sistema nunca substitua o julgamento clínico humano.

Este projeto foi desenvolvido no contexto do \textbf{Tech Challenge da Fase~3} do curso de Pós-Graduação em Inteligência Artificial para Devs da FIAP, cujo tema central é \textbf{IA Generativa}. O desafio propõe a construção de um assistente virtual inteligente que integre técnicas de fine-tuning de LLMs, orquestração multi-agente com LangGraph e segurança por design.

\section{Objetivos do Projeto}

O objetivo principal deste trabalho é desenvolver um \textbf{Assistente Virtual Médico} capaz de auxiliar profissionais de saúde no acesso a informações clínicas contextualizadas. Para isso, o sistema integra cinco pilares fundamentais:

\begin{enumerate}
    \item \textbf{Fine-tuning} do modelo Mistral 7B Instruct v0.3 com o dataset MedQuAD traduzido para português, utilizando \acrfull{qlora} 4-bit no Google Colab;
    \item \textbf{\acrfull{rag}} com 14 Linhas de Cuidado reais (protocolos clínicos profissionais) indexadas no ChromaDB;
    \item \textbf{Grafo multi-agente} implementado em LangGraph com 6 agentes especializados;
    \item \textbf{Segurança} com guardrails que impedem prescrições diretas e exigem validação humana;
    \item \textbf{Interface} Gradio para demonstração interativa do assistente.
\end{enumerate}

\begin{importantbox}[Aviso de Segurança]
Este sistema foi projetado exclusivamente como ferramenta de \textbf{suporte à decisão clínica}, nunca como substituto do profissional de saúde. Todas as respostas incluem avisos de validação humana obrigatória e o sistema é impedido, por guardrails, de prescrever medicamentos diretamente.
\end{importantbox}

\section{Arquitetura Geral do Sistema}

A Figura~\ref{fig:arquitetura_geral} apresenta a visão de alto nível da arquitetura, organizada em cinco camadas que correspondem aos cinco pilares do projeto.

\begin{figure}[H]
\centering
\begin{tikzpicture}[
    node distance=0.8cm and 1.2cm,
    block/.style={rectangle, draw, rounded corners=3pt, minimum width=3.8cm, minimum height=1cm, align=center, font=\small\sffamily},
    pillar/.style={block, fill=fiap-red!15, draw=fiap-red, thick},
    data/.style={block, fill=blue!10, draw=blue!60, thick},
    agent/.style={block, fill=green!10, draw=green!60!black, thick},
    io/.style={block, fill=orange!10, draw=orange!70, thick},
    arrow/.style={->, thick, >=stealth}
]

% Camada de Dados
\node[data] (medquad) {MedQuAD\\(16.400 pares QA)};
\node[data, right=of medquad] (linhas) {14 Linhas de\\Cuidado (PDF)};
\node[data, right=of linhas] (sinteticos) {Dados Sintéticos\\(80 pacientes)};

% Camada de Processamento
\node[pillar, below=1.2cm of medquad] (finetune) {\textbf{Fine-Tuning}\\QLoRA 4-bit};
\node[pillar, below=1.2cm of linhas] (rag) {\textbf{RAG}\\ChromaDB};
\node[pillar, below=1.2cm of sinteticos] (postgres) {\textbf{PostgreSQL}\\Dados Clínicos};

% Camada Multi-Agente (centralizada)
\node[agent, below=1.2cm of rag, minimum width=12.6cm] (grafo) {
    \textbf{Grafo Multi-Agente LangGraph} --- 6 Agentes\\
    Triagem $\rightarrow$ Protocolo / Paciente Data $\rightarrow$ Raciocínio $\rightarrow$ Guardrails $\rightarrow$ Explicabilidade
};

% Camada de Segurança
\node[pillar, below=1.2cm of grafo, minimum width=12.6cm] (seguranca) {
    \textbf{Segurança}: Guardrails + Logging de Auditoria + Explicabilidade
};

% Camada de Interface
\node[io, below=0.8cm of seguranca] (gradio) {\textbf{Interface Gradio}};

% Setas
\draw[arrow] (medquad) -- (finetune);
\draw[arrow] (linhas) -- (rag);
\draw[arrow] (sinteticos) -- (postgres);
\draw[arrow] (finetune) -- (grafo.north -| finetune);
\draw[arrow] (rag) -- (grafo);
\draw[arrow] (postgres) -- (grafo.north -| postgres);
\draw[arrow] (grafo) -- (seguranca);
\draw[arrow] (seguranca) -- (gradio);

\end{tikzpicture}
\caption{Arquitetura geral do Assistente Virtual Médico.}
\label{fig:arquitetura_geral}
\end{figure}

\section{Guia de Reprodução Passo a Passo}\label{sec:reproducao}

O projeto foi concebido para ser \textbf{integralmente reprodutível}. Um pipeline de 7 scripts numerados, localizado no diretório \texttt{scripts/}, permite ao avaliador reconstruir todo o sistema do zero, desde a geração de dados sintéticos até a execução do assistente completo.

\subsection{Pipeline de Scripts}

A Tabela~\ref{tab:pipeline_scripts} descreve cada script, suas pré-condições e saídas.

\begin{table}[H]
\centering
\caption{Pipeline de reprodução: scripts numerados.}
\label{tab:pipeline_scripts}
\rowcolors{2}{fiap-light}{white}
\begin{tabularx}{\textwidth}{c l X X}
\toprule
\textbf{\#} & \textbf{Script} & \textbf{O que faz} & \textbf{Saída} \\
\midrule
01 & \texttt{gerar\_dados\_sinteticos} & Gera pacientes fictícios com comorbidades, exames, prontuários e receitas clinicamente coerentes & \texttt{data/synthetic/*.json} (4 arquivos) \\
02 & \texttt{setup\_postgres} & Cria container Docker com PostgreSQL, tabelas e popula com dados do script 01 & Banco \texttt{medical\_assistant\_db} operacional \\
03 & \texttt{preparar\_dataset} & Baixa MedQuAD do HuggingFace, traduz para PT-BR e formata em Alpaca & \texttt{data/processed/training\_data.jsonl} \\
\rowcolor{orange!10}
-- & \textit{Pausa: Google Colab} & \textit{Fine-tuning QLoRA no notebook} \texttt{03\_fine\_tuning.ipynb} & \textit{Modelo em} \texttt{models/} \\
04 & \texttt{indexar\_protocolos} & Chunking dos protocolos Markdown e indexação no ChromaDB & Banco vetorial \texttt{chroma\_db/} \\
05 & \texttt{avaliar\_baseline} & Roda 50 perguntas médicas no Mistral 7B base & \texttt{data/evaluation/baseline.json} \\
06 & \texttt{avaliar\_finetuned} & Mesmas 50 perguntas no modelo fine-tuned, com comparação e métricas & Relatório comparativo \\
07 & \texttt{iniciar\_app} & Inicia o assistente completo com interface Gradio & Servidor web local \\
\bottomrule
\end{tabularx}
\end{table}

\subsection{Fluxo de Execução}

\begin{figure}[H]
\centering
\begin{tikzpicture}[
    node distance=0.4cm,
    pstep/.style={rectangle, draw=fiap-dark, rounded corners=3pt, fill=fiap-light, minimum width=2.2cm, minimum height=0.8cm, align=center, font=\small\sffamily\bfseries},
    colab/.style={pstep, fill=orange!20, draw=orange!70},
    arrow/.style={->, thick, >=stealth, color=fiap-red}
]
    \node[pstep] (s01) {01};
    \node[pstep, right=of s01] (s02) {02};
    \node[pstep, right=of s02] (s03) {03};
    \node[colab, right=of s03] (colab) {Colab};
    \node[pstep, right=of colab] (s04) {04};
    \node[pstep, right=of s04] (s05) {05};
    \node[pstep, right=of s05] (s06) {06};
    \node[pstep, right=of s06] (s07) {07};

    \draw[arrow] (s01) -- (s02);
    \draw[arrow] (s02) -- (s03);
    \draw[arrow] (s03) -- (colab);
    \draw[arrow] (colab) -- (s04);
    \draw[arrow] (s04) -- (s05);
    \draw[arrow] (s05) -- (s06);
    \draw[arrow] (s06) -- (s07);

    % Labels abaixo
    \node[below=0.3cm of s01, font=\tiny\sffamily, text width=2cm, align=center] {Dados\\Sintéticos};
    \node[below=0.3cm of s02, font=\tiny\sffamily, text width=2cm, align=center] {PostgreSQL};
    \node[below=0.3cm of s03, font=\tiny\sffamily, text width=2cm, align=center] {Dataset\\MedQuAD};
    \node[below=0.3cm of colab, font=\tiny\sffamily, text width=2cm, align=center] {Fine-Tuning\\QLoRA};
    \node[below=0.3cm of s04, font=\tiny\sffamily, text width=2cm, align=center] {RAG\\ChromaDB};
    \node[below=0.3cm of s05, font=\tiny\sffamily, text width=2cm, align=center] {Baseline};
    \node[below=0.3cm of s06, font=\tiny\sffamily, text width=2cm, align=center] {Avaliação\\Fine-Tuned};
    \node[below=0.3cm of s07, font=\tiny\sffamily, text width=2cm, align=center] {Assistente\\Gradio};
\end{tikzpicture}
\caption{Fluxo de execução do pipeline de reprodução.}
\label{fig:fluxo_pipeline}
\end{figure}

\subsection{Como Reproduzir}

\begin{lstlisting}[style=bash, caption={Comandos para reprodução completa do projeto.}]
# 0. Pre-requisitos: Python 3.11+, Docker Desktop
git clone https://github.com/Zagari/virtual-medical-assistant.git
cd virtual-medical-assistant
pip install -r requirements.txt

# 1. Gerar dados sinteticos (80 pacientes ficticios)
python scripts/01_gerar_dados_sinteticos.py

# 2. Subir PostgreSQL e popular banco
python scripts/02_setup_postgres.py

# 3. Preparar dataset MedQuAD para fine-tuning
python scripts/03_preparar_dataset.py

# --- PAUSA: Fine-tuning no Google Colab ---
# Abrir notebooks/03_fine_tuning.ipynb no Colab
# Upload de data/processed/training_data.jsonl
# Treinar modelo -> Download para models/

# 4. Indexar protocolos clinicos no ChromaDB
python scripts/04_indexar_protocolos.py

# 5-6. Avaliar modelo (baseline vs fine-tuned)
python scripts/05_avaliar_baseline.py
python scripts/06_avaliar_finetuned.py

# 7. Iniciar o assistente
python scripts/07_iniciar_app.py
\end{lstlisting}

\begin{infobox}[Princípios do Pipeline]
Cada script é \textbf{autocontido} e pode ser executado independentemente. Todos imprimem progresso durante a execução (registros gerados, tempo decorrido) e \textbf{validam pré-condições} antes de iniciar (por exemplo, o script~02 verifica se os arquivos de dados sintéticos existem antes de tentar popular o banco).
\end{infobox}

\section{Estrutura de Arquivos do Projeto}

\begin{lstlisting}[style=bash, caption={Estrutura de diretórios do projeto.}, label=lst:estrutura]
virtual-medical-assistant/
|-- README.md
|-- requirements.txt
|-- .env.example
|-- scripts/                    # Pipeline numerado (01 a 07)
|-- data/
|   |-- raw/                    # MedQuAD original
|   |-- processed/              # Dataset formatado (training_data.jsonl)
|   |-- linhas_de_cuidado/      # 14 PDFs + extracted/ (Markdown)
|   +-- synthetic/              # Pacientes, exames, prontuarios, receitas
|-- notebooks/
|   +-- 03_fine_tuning.ipynb    # Notebook para Google Colab
|-- src/
|   |-- data/                   # Geracao e preprocessamento
|   |-- fine_tuning/            # Configuracao e treino
|   |-- assistant/              # LLM wrapper, RAG, prompts
|   |-- flows/                  # LangGraph (state, graph, agents/)
|   |-- security/               # Guardrails, logging, auditoria
|   +-- database/               # SQLAlchemy models e queries
|-- app/                        # Interface Gradio
|-- docs/                       # Relatorio LaTeX
+-- tests/                      # Testes automatizados
\end{lstlisting}

\section{Organização do Documento}

O restante deste relatório está organizado da seguinte forma:

\begin{itemize}
    \item \textbf{Capítulo~\ref{cap:dados}} apresenta os dados utilizados: as 14 Linhas de Cuidado, a geração de dados sintéticos, o banco de dados PostgreSQL e o dataset MedQuAD;
    \item \textbf{Capítulo~3} detalha o processo de fine-tuning do Mistral 7B com QLoRA;
    \item \textbf{Capítulo~4} descreve a implementação do RAG com ChromaDB;
    \item \textbf{Capítulo~5} apresenta o grafo multi-agente com LangGraph;
    \item \textbf{Capítulo~6} aborda a segurança, logging e interface Gradio;
    \item \textbf{Capítulo~7} traz a avaliação e os resultados obtidos;
    \item \textbf{Capítulo~8} conclui com reflexões, limitações e trabalhos futuros.
\end{itemize}
